% ------------------------------------------------------------------------------
% layout
\documentclass [twoside]{scrreprt}                % [] 10pt|11pt|12pt font size, default 10pt
	  \KOMAoptions{paper=a4}

% encoding stuff
\usepackage [utf8]{inputenc}
\usepackage [T1]{fontenc}                         % fuer korrekte trennung deutscher umlaute, scharf-s etc.


% figures
\usepackage {pgf}                                 % includepgf for bitmaps
\usepackage {epsfig}
\usepackage {psfrag}                              % replacing strings in eps files
%\usepackage [bf,small]{caption}                   % caption style
\usepackage {ifpdf}
%\ifpdf
    \usepackage{graphicx}
    \usepackage{epstopdf}
    %\DeclareGraphicsRule{.eps}{pdf}{.pdf}{`epstopdf #1}
    \pdfcompresslevel=9
%\else
  %  \usepackage {graphicx}
%\fi

\usepackage {subfig}                              % subfigures a) b)...
%\usepackage {floatflt}                           % packages supporting in-text floats
%\usepackage {wrapfig}                            % another one
%\usepackage {picins}                             % another one
%\usepackage{timing}                              % deprecated use, use below!
\usepackage{sidecap}

\usepackage[overlay]{textpos}
\usepackage{listings}
\lstset{
    frameround=fttt,
    language=C,
    %numbers=left,
    breaklines=true,
    keywordstyle=\color{red}\bfseries, 
    basicstyle=\ttfamily\scriptsize,
    numberstyle=\color{black}
    }
\lstMakeShortInline[columns=fixed]|
\usepackage{multicol}
\usepackage{tcolorbox}
\tcbuselibrary{skins,breakable}
\usepackage{tikz}
\usetikzlibrary{shadings,shadows}
\newenvironment{myexampleblock}[1]{%
    \tcolorbox[beamer,%
    noparskip,breakable,%
    colbacklower=LimeGreen!75!LightGreen,%
    title=#1]}%
    {\endtcolorbox}
%\newenvironment<>{varblock}[2][.9\textwidth]{%
%  \setlength{\textwidth}{#1}
%  \begin{actionenv}#3%
%    \def\insertblocktitle{#2}%
%    \par%    
%    \usebeamertemplate{myexampleblock begin}%
%  }
%  {\par%
%  \usebeamertemplate{myexampleblock end}%
%  \end{actionenv}}
%
%\newenvironment{myblock}[1]{\begin{textblock*}{500pt}(#1)}{\end{textblock*}}
%
%% special way to reach the top of the block
%\def\newabove(#1){
%  ([yshift=1.5ex]#1.center)
%}

% tables
\usepackage{multirow}                             % col- and rowspan...
\usepackage{booktabs}                             % professional table layout


% math
\usepackage{bbm}                                  % black board style, (e.g. \mathbbm{R})
\usepackage{amssymb}
\usepackage{amsmath}                              % e.g. \nobreakslash- for non-breakable hyphens
\usepackage{bbold}
\usepackage{amsfonts}
\usepackage{units}

% text
\usepackage {latexsym}                            % additional symbols
%\usepackage{lmodern}                              % latin modern
\usepackage [english,ngerman]{babel}              % language (change via "\selectlanguage{english}")
%\usepackage[colorlinks=true,pdftex=true,linkcolor=black,citecolor=black,urlcolor=blue,plainpages=false,pdfpagelabels]{hyperref}
\usepackage[colorlinks=true,linkcolor=black,citecolor=black,filecolor=black,urlcolor=black,plainpages=false,pdfpagelabels,breaklinks=true]{hyperref}
%\usepackage[all]{hypcap}                          % links to figures point to upper border (not caption)
\usepackage{breakurl}
\usepackage{hyphenat}
\usepackage[normalem]{ulem}                       % underline text...
\usepackage{listings}                            % source/code listings
\usepackage[plain]{algorithm}
\usepackage{algorithmicx}                        % pseudo code
\usepackage{algpseudocode}                       % pseudo code
\usepackage{units}                               % \unit[123]{U}
\usepackage{mdwlist}                             % dichtere itemize* etc.
\usepackage{paralist}                            % inline Listen
\usepackage{natbib}


% list of abbreviations, index
\usepackage{nomencl}
    \let\abbrev\nomenclature
    \renewcommand{\nomname}{List of Abbreviations}
    \setlength{\nomlabelwidth}{.25\hsize}
    \renewcommand{\nomlabel}[1]{#1 \dotfill}
    \setlength{\nomitemsep}{-\parsep}
    %\makeglossary
    \makenomenclature
    \newcommand{\markup}[1]{\uline{#1}}
\usepackage{makeidx}                              % index of key words
    \makeindex


% layout
    \usepackage[colorinlistoftodos,prependcaption,textsize=tiny]{todonotes}
%\usepackage{stfloats}
%\usepackage{dblfloatfix}
\usepackage{fixltx2e}
%\usepackage {afterpage}                          % \afterpage{xyz} executes xyz after current page is finished
%\usepackage {flafter}                            % no object before definition
%\usepackage {float}                              % Here(!!!) placement option
% ------------------------------------------------------------------------------


% source aliases
\newcommand{\thesisTitle}{Bachelor Thesis}
\newcommand{\thesisTitleGerman}{German Title}
\newcommand{\thesisAuthor}{Arthur Heimbrecht}
\newcommand{\thesisAuthorBornIn}{Speyer}



% hypenation (englisch)
\selectlanguage{english}
\hyphenation    {HAGEN HASTE HANNEE Bert-schin-ger Maass ma-xi-mum SoftHASTE
neu-ron pre-synaptic post-synaptic}

% hypenation (deutsch)
\selectlanguage{ngerman}
\hyphenation    {HAGEN HASTE HANNEE Bert-schin-ger Maass SoftHASTE}


\tolerance 1414
\hbadness 1414
\emergencystretch 1.5em % critical parameter!!!
%\displaywidowpenalty % avoid widows after math displays
\raggedbottom % should be default anyway
%\clubpenalty = 10000
%\widowpenalty = 10000
%\hfuzz 0.3pt
%\vfuzz

\newcommand{\tikzmark}[2][minimum width=6cm,minimum height=1.5cm]{
	\tikz[remember picture,overlay]
	\node[anchor=west,
	      inner sep=0pt,
	      outer sep=6pt,
	      xshift=-0.5em,
	      yshift=-3ex,
	      #1](#2){};
}

\newcommand{\shownode}[1]{
       \tikz[remember picture,overlay]\draw[red](#1.south east)rectangle(#1.north west);
}

\newcommand{\showanchor}[1]{
       \tikz[remember picture,overlay]\draw[red,thick,mark=x] plot coordinates{(#1)};
}
%\renewcommand{\bottomfraction} {0.70} % max fraction of page for floats at bottom
%\renewcommand{\textfraction} {0.20} % min fraction of page for text
%\renewcommand{\floatpagefraction} {0.60} % min fraction of floatpage that should have floats
%\renewcommand{\dblfloatpagefraction}{0.60} % for twocolumn layout, see floatpagefraction above
%\setcounter{totalnumber} {6} % max number of floats on a single page
%\setcounter{topnumber} {6} % max number of top floats on a single page
%\setcounter{dbltopnumber}{6} % for twocolumn layout, see topnumber above
%\setcounter{bottomnumber}{6} % max number of bottom floats on a single page
% ------------------------------------------------------------------------------

\pagestyle  {empty}                               % plain, empty, headings oder myheadings.



\begin{document}

\tikzstyle{every picture}+=[remember picture]

% ------------- Cover Pages ----------------------

\begin{titlepage}

\textwidth 		14cm
\textheight 		27.0cm

\evensidemargin    	-1.5cm
\oddsidemargin   	-1.5cm
\topmargin     		-3.0cm

\thispagestyle{empty}


\begin{figure}
\includegraphics{fig/cover/Publ-Cover1.eps}
%\epsfig{file=./fig/cover/Publ_Cover1.eps}% , width=\textwidth}%,height=4.5cm}
\end{figure}

\vspace*{3cm}

\hspace{5cm}
\begin{minipage}{8.2cm}
\begin{flushright}
\Large{\sf{\thesisAuthor}}
\end{flushright}
\end{minipage}


\begin{figure}[ h]
\includegraphics{fig/cover/Publ-Cover2.eps}
%\epsfig{file=./fig/cover/Publ_Cover2.eps}%,width=12.7cm}
\end{figure}

\vspace*{3mm}

\hspace{3cm}
\begin{minipage}{10.2cm}
\begin{flushright}
\Large{\sf{\thesisTitle}}
\end{flushright}

\vspace*{3mm}

\vspace*{4mm}

\begin{flushright} \Large{\sf{
%%% no = fortlaufende KIP-Nummer, wird von der Bibliothek (Tel. 8992) vergeben.
%HD-KIP-04-16
}}
\end{flushright}

\end{minipage}


\begin{figure} [b]
\includegraphics{fig/cover/Publ-Cover3.eps}
%\epsfig{file=./fig/cover/Publ_Cover3.eps}%,width=12.7cm}
\end{figure}

\end{titlepage}


%-------------------------------------------------






% and now in englisch!
\selectlanguage{english}





% ------------- Abstract -------------------------


\pagestyle{plain}

\setcounter{page}{1}
\pagenumbering{Roman}

\begin{abstract}
\begin{center}
\textbf{\thesisTitle}
\end{center}

As part of the Human Brain Project, BrainScaleS is a unique project on many levels. This includes a processor solely used by the HICANN-DLS, which manages synaptic weights for every neuron built into one of the many wafers. To accelerate the speed at which this so called plasticity processor unit (PPU) computes all synaptic weights of every neuron used, the processor has an extended instruction set architecture (ISA) that supports vector registers and single input multiple data (SIMD).
This report deals with the task of adding built-in functions to an existing back-end of GCC, specifically the one used by the PPU, in order to extend the already implemented set of functions according to the users needs.
\end{abstract}


%-------------------------------------------------




% ------------- Table of Contents ----------------

\cleardoublepage
\setcounter{tocdepth}{3}
\tableofcontents

%-------------------------------------------------






% ------------- Chapters -------------------------

\cleardoublepage

\pagestyle{headings}

\setcounter{page}{1}
\pagenumbering{arabic}



%%% INCLUDES %%%%%%%%%%%%%%%%%%%%%%%%%%
% introduction
\chapter{Compiler Structure}
\label{chapter:compiler structure}
\newlength\Colsep
\setlength\Colsep{10pt}
Typically a compiler consists of a front-end, middle-end and back-end. These three parts sit on top of each other with the front-end on top and the back-end at the bottom and pass down the programs code as it is translated and optimized or compiled. But communication between the parts does not go only one way (at least for GCC)\todo{add reference}
and changes that are made in the back-end affect the front-end as well!
Such modular architecture  allows for a compiler to support many different front-ends and back-ends independently because all code needs to pass the same middle-end no matter which front-end or back-end is used. In normal cases a front-ends are associated with programming languages and back-ends with CPU architectures or ``targets'' while the middle-end is at the heart of the compiler and is usually what distinguishes between different compilers, as does the communication between the three parts.
\smallskip
\noindent\begin{minipage}{\textwidth}
\begin{minipage}[c]{\dimexpr0.2\textwidth-0.5\Colsep\relax}
    \begin{tikzpicture}[baseline=(current bounding box.south west)]
    	\coordinate (anch) at (0,0);
    	%%\coordinate (anch2) at (0,13);
		\node (tl) at (2.8, 10) {\texttt{o}};
		\node (ml) at (2.7, 6) {\texttt{o}};
		\node (bl) at (2.7, 3) {\texttt{o}};
		%%\draw[->] (anch) -- (tl);
	\end{tikzpicture}
\end{minipage}\hfill
\begin{minipage}[c]{\dimexpr0.6\textwidth-0.5\Colsep\relax}
 \begin{myexampleblock}{front-end}
	\begin{itemize}
		\item recognizes programming language \tikz[baseline] \node[coordinate] (box1){}; 
		\item pre-processing
		\item type checking
		\item generates IR
	\end{itemize}
\end{myexampleblock}
\begin{myexampleblock}{middle-end}
	\begin{itemize}
		\item general optimizations to IR \tikz[baseline] \node[coordinate] (top2) {};
		\item compiler specific actions
		\item passes IR \tikz[baseline] \node[coordinate] (bot2) {};
	\end{itemize}
\end{myexampleblock}
\begin{myexampleblock}{back-end}
	\begin{itemize}
		\item target specific optimization \tikz[baseline] \node[coordinate] (top3) {};
		\item register allocation (spilling)
		\item generation of assembly code
	\end{itemize}
\end{myexampleblock}
\end{minipage}
\begin{minipage}[c]{\dimexpr0.2\textwidth-0.5\Colsep\relax}
    \begin{tikzpicture}
		\node at (0, 10.5) {\texttt{o}};
		\node at (0, 6.5) {\texttt{o}};
		\node at (0, 3) {\texttt{o}};
	\end{tikzpicture}
\end{minipage}\hfill
\end{minipage}
\vspace{\baselineskip}
The first part of the compilation process is the translation of code which is written in some programming language by the front-end into a so called Intermediate Representation (IR) that looks the same for every front-end language and usually is never seen by the user. AS already described, any supported programming language (C, C++, Java…) is implemented in its own front-end that defines how the language is translated into IR. After that the IR is send to the middle-end where it is mainly optimized and then passed to the back-end. The back-end first performs further optimization that is target-specific followed by allocating registers and handling relative memory. Finally the code is translated into the assembly language that is supported by the target which the back-end belongs to.
This last step executed by the back-end can seen as independent from the back-end as the translation into assembly is pretty straight forward, but the rs/6000 back-end for instance allows for direct assembly output that is not dependent on IR. Therefore we combine these the assembler and the back-end.
After the code is compiled and emitted as an object file it is also linked, which means combining different object files and assigning absolute memory addresses to them and at last the symbols which are used in assembly are substituted by their respective opcodes. Opcodes are almost the lowest representation of machine code as the machine would not accept strings a input but rather bits that are represented as hex literals. These opcodes form pairs with assembly macros and can be found in binutils. Finally the resulting binary file is emitted by the linker and loaded into the memory of the processor to be executed.

GCC generally obeys this scheme and therefore the PowerPC architecture (also called rs/6000) is a good example of a back-end that supports different extensions to the standard architecture such as the AltiVec vector extension. This makes is more complex than the Cell Synergistic Processor Units (SPU) back-end but all back-ends in GCC share the same functions defined by macros. Still besides those macros the back-ends obey only few restrictions by the middle-end and thus this guide will not easily be transferable to other GCC back-ends and especially not to other compilers back-ends!

% chapter X
\chapter{Creating an intrinsic function}
\label{chapter:builtins}

An intrinsic or builtin function is a medium link between inline assembly and normal functions. They look like normal functions but usually trigger certain machine instructions on a very basic level. This makes intrinsic functions highly effective compared to normal functions. Intrinsic functions are usually directly built into the compiler or rather its back-end and provide the compiler with additional information that allows for highly optimized code. Usually there is one specific machine instruction at the core of every intrinsic function.
Most of this is done in the machine description that builds an interface between machine instructions and IR. Every instruction that is used within the compiler is part of a so called insn (short for instruction). An insn is an expression that is identified by its code and adds information to a machine instruction. It is written in Register Transfer Language (RTL) that has its own syntax and keywords.
We will not dive further into this specific topic since these details are not needed for the purpose of this report. Therefore we will utilize existing insns and see what makes them complete intrinsic functions.

The functions we will take a look at is called |vec_addc| or |vec_vaddcuw| when used in any front-end. This is one of the differences between normal functions and intrinsics as the latter one is the same for every front-end because it is implemented into the back-end.
This function, which takes two vectors that have unsigned int elements as arguments, "returns a vector containing the carries produced by adding each set of corresponding elements of two given vectors" , this means for |c = vec_addc(a,b)| that the resulting vector |c|'s elements' bits are 1 if adding |a| and |b| produces a carry at that bit position and |0| otherwise. Therefore the name "Vector ADD Carry Unsigned Wordsize". But the use of this function is less important than the way it is implemented.\todo{https://www.ibm.com/ support/knowledgecenter/ SSGH2K\_13.1.3/com.ibm.xlc1313. aix.doc/compiler\_ref/ vec\_addc.html}
The alternative acronym vaddcuw depicts the name of the instruction on a processor level which is "|vaddcuw  %c,%a,%b|" and is implemented as a synonym intrinsic to |vec_addc|. Hence |vaddcuw| and |vec_addc| have their specific insn with the code |altivec_vaddcuw|.
\begin{lstlisting}[language=C++,basicstyle=\ttfamily\scriptsize,keywordstyle=\color{red}]
(define_insn "altivec_vaddcuw"
[(set
	(match_operand:V4SI 0 "register_operand" "=v")
	(unspec:V4SI [(match_operand:V4SI 1 "register_operand" "v")
			(match_operand:V4SI 2 "register_operand" "v")] UNSPEC_VADDCUW))]
"VECTOR_UNIT_ALTIVEC_P (V4SImode)"
"vaddcuw %0,%1,%2"
[(set_attr "type" "vecsimple")])
\end{lstlisting}
To explain this in a few words: the first line depicts the insn code |altivec_vaddcuw|, the next four lines specify what the instruction does and which constriants the operands (0,1 and 2) need to fulfill (this is RTL). It follows a boolean function, to check if this insn is valid to be used and the instructions name with its operands. The attributes set at the end are only used internally.

Now our first step is to add an entry into the builtin-description-file rs6000-builtin.def. This file holds all builtin functions and connects them to insns. The very beginning of this file consists of convenience macros for different extensions that allow for better readability as most of the properties of a builtin function are similar or even the same for each extension.
\begin{lstlisting}[language=C++,basicstyle=\ttfamily\scriptsize,keywordstyle=\color{red}]
#define BU_ALTIVEC_1(ENUM, NAME, ATTR, ICODE)				\
  RS6000_BUILTIN_1 (ALTIVEC_BUILTIN_ ## ENUM,		/* ENUM */	\
  		    "__builtin_altivec_" NAME,		/* NAME */	\
  		    RS6000_BTM_ALTIVEC,			/* MASK */	\
  		    (RS6000_BTC_ ## ATTR		/* ATTR */	\
   		     | RS6000_BTC_UNARY),				\
  		    CODE_FOR_ ## ICODE)			/* ICODE */
\end{lstlisting}
There are different types of macros for different types of builtin functions. Each macro takes four arguments besides it's enumeration name which are: The name of the function as string literal, a bit-mask that indicates which options are enabled, attribute information and the insn code. \todo{make this an actual list}
The macros then add information for each extension and divide the builtin functions into certain groups which depend on the number of arguments. In our case the macros are called |BU_ALTIVEC_1, BU_ALTIVEC_2, BU_ALTIVEC_3| for up to three arguments respectively and all of these have the |RS6000_BTM_ALTIVEC| bit-mask as well as the same prefix for their enumeration name and function name. These are |ALTIVEC_BUILTIN_| for the enumeration name and |__builtin_altivec_| for the function name. Besides that each macro has a specific attribute such as |RS6000_BTC_TERNARY| for function with three arguments.

Now we finally move to specifying the builtin function. The line of code we are interested in is:
\begin{lstlisting}[language=C++,basicstyle=\ttfamily\scriptsize,keywordstyle=\color{red}]
BU_ALTIVEC_2 (VADDCUW, "vaddcuw", CONST, altivec_vaddcuw)
\end{lstlisting}
First the name for the enumeration is set as |ALTIVEC_BUILTIN_VADDCUW| then the builtin functions name is set as |__builtin_altivec_vaddcuw|. Next the attribute is set as |CONST| which means that no other registers are altered when the insn is used but the three registers that are directly called. At last the insn code is given as |altivec_vaddcuw|, which we know from earlier. This already gives us a usable builtin function! To use it we first set our vector variables:
\begin{lstlisting}[language=C++,basicstyle=\ttfamily\scriptsize,keywordstyle=\color{red}]
vector unsigned int a,b,c;
c = __builtin_altivec_vaddcuw(a, b);
\end{lstlisting}
But this function still has some flaws as it would not give an error for this case:
\begin{lstlisting}[language=C++,basicstyle=\ttfamily\scriptsize,keywordstyle=\color{red}]
short a,b,c;
c = __builtin_altivec_vaddcuw(a, b);
\end{lstlisting}
Because there is no type-checking a user could use the function in a completely wrong manner. To avoid this though, there are overloaded builtin functions that include a type-checking routine.

An overloaded builtin function is basically just another builtin function that is less specific than the previous function as it only specifies two names and no insn code. Thus they have their own macros (|BU_ALTIVEC_OVERLOADED_1, BU_ALTIVEC_OVERLOADED_2, BU_ALTIVEC_OVERLOADED_3|) and differ in a way that the enumeration prefix is |ALTIVEC_BUILTIN_VEC_| and the function name prefix is |__builtin_vec_| also the attributes are completely set in advance as well as the specific attributes.
For an overloaded builtin we will use a simpler name which will be  |__builtin_vec_addc|
\begin{lstlisting}[language=C++,basicstyle=\ttfamily\scriptsize,keywordstyle=\color{red}]
BU_ALTIVEC_OVERLOAD_2 (ADDC, "addc")
\end{lstlisting}
Since there is no insn code given all overloaded builtins are connected to existing builtins which we will do next in rs6000-c.c.

We will overload the builtin |__builtin_vec_addc| with |__builtin_altivec_vaddcuw| and add argument and return types. In principle this is similar to a functions argument types but for intrinsics these are declared in a |struct| that allows for different combinations of argument and return types. The |struct| is called |altivec_builtin_types| and consists of an overloaded builtin code, a normal builtin code, the return type and up to three argument types. For |ADDC| exists only one |struct| though because it only works for vectors of unsigned |ints| (|V4SI| = Vector of 4 Single Integers):
\begin{lstlisting}[language=C++,basicstyle=\ttfamily\scriptsize,keywordstyle=\color{red}]
{ ALTIVEC_BUILTIN_VEC_ADDC, ALTIVEC_BUILTIN_VADDCUW,
RS6000_BTI_unsigned_V4SI, RS6000_BTI_unsigned_V4SI, RS6000_BTI_unsigned_V4SI, 0 }			
\end{lstlisting}
|ALTIVEC_BUILTIN_VEC_ADDC| is now overloaded with the working builtin function |ALTIVEC_BUILTIN_VADDCUW| from earlier and has a return type and argument types which are vectors of unsigned |ints|. The last entry is 0 because there is no third argument.

Basically this is enough for the overloaded builtin function to work properly and it can be used in a way such as
\begin{lstlisting}[language=C++,basicstyle=\ttfamily\scriptsize,keywordstyle=\color{red}] 
vector unsigned int a,b,c;
c = __builtin_vec_addc(a, b);
\end{lstlisting}
and would give an error if the types would not match those we set earlier.

To give all of this a nice touch and increase usability in the end. We define synonyms for our newly created overloaded builtin function. We will not do this for the original builtin function since we want to avoid the missing type checking.
\begin{lstlisting}[language=C++,basicstyle=\ttfamily\scriptsize,keywordstyle=\color{red}] 
#define vec_vaddcuw vec_addc
...
#define vec_addc __builtin_vec_addc
\end{lstlisting}
Here the first line defines a synonym for the function for people familiar with the assembly macro.
This brings our task of defining a builtin function to an end!

But there is still one kind of common builtin function left that differs to normal one-to-three-argument-builtins in some ways such as requiring a memory address for assembly macro instead of a register or simply not having a return value. These special builtin functions that have the convenience macros |BU_ALTIVEC_X| and |BU_ALTIVEC_OVERLOADED_X| though a special macro can also be overloaded with a normal overload macro like |BU_ALTIVEC_OVERLOADED_2|.
The special X-macro has |CODE_FOR_nothing| as insn code like the overloaded macros did earlier and thus is not intended to be handled normally but will be caught in the main file |rs6000.c| which we will see later.

We will have an example that uses a normal overloaded macro since it is slightly easier and special overloaded functions tend to need special handling in the main back-end file because the have very specific properties. Therefore we take a look at |vec_mtvscr(a)| which copies the value of |a| into the Vector status and Control Register (|VSCR|) (Move To |VSCR|). The insn code for this builtin function is |altivec_mtvscr| and the machine instruction is |mtvscr %a|.
It is obvious that this function does not generate any return value and therefore would not fit a one-argument-builtin.
\begin{lstlisting}[language=C++,basicstyle=\ttfamily\scriptsize,keywordstyle=\color{red}] 
BU_ALTIVEC_X (MTVSCR, "mtvscr",	MISC)
...
BU_ALTIVEC_OVERLOAD_1 (MTVSCR, "mtvscr")
\end{lstlisting}
The other difference for this builtin is that it is not a |CONST| builtin but carries a |MISC| attribute. This argument is only used in special cases that make an exception to |CONST| or any of the other special attributes and means specifically that there are no special attributes. We will not discuss the other attributes but an explanation can be found in the |rs6000.h| file.
In contrast to the builtin function |__biultin_altivec_vaddcuw| from earlier, the builtin function |__builtin_altivec_mtvscr| is of no use now since there is no insn code connected with it. Thus we will add special cases in the function that handles the builtin functions or "expands" them. This is done in the |altivec_expand_builtin| function that handles special builtins exclusively. Normal builtins are expanded depending on their number of arguments at the very end of |rs6000_expand_builtin| where |altivec_expand_bulitin| is called before hand. 

In the expander function the compiler switches between all special cases, which means there has to be an entry for every special builtin there is. For |mtvsrc| this entry looks something like:
\begin{lstlisting}[language=C++,basicstyle=\ttfamily\scriptsize,keywordstyle=\color{red}]
case ALTIVEC_BUILTIN_MTVSCR:
      icode = CODE_FOR_altivec_mtvscr;
      arg0 = CALL_EXPR_ARG (exp, 0);
      op0 = expand_normal (arg0);
      mode0 = insn_data[icode].operand[0].mode;

      /* If we got invalid arguments bail out before generating bad rtl.  */
      if (arg0 == error_mark_node)
    	return const0_rtx;

      if (! (*insn_data[icode].operand[0].predicate) (op0, mode0))
     	op0 = copy_to_mode_reg (mode0, op0);

      pat = GEN_FCN (icode) (op0);
      if (pat)
    	emit_insn (pat);
      return NULL_RTX;
\end{lstlisting}
This code is probably the most difficult part when adding a special builtin function. The easiest way is to look for a similar function, copy its code and modify it if necessary. But we will go through this code briefly:
First we see some important variables that get their respective values. |icode| obviously holds the insn code, |arg0| holds whatever the function gets as first argument, |op0| gets the operand of that argument, and |mode| holds the mode of the operand that the insn needs. It then checks if the argument is actually valid and returns an error otherwise. Next it checks whether |mode| and |op0| match and tries to convert the operand if they do not match. |pat| gets to hold the directive to build an insn with code |icode| and operand |op0| and if this gives no error the final insn is emitted. The return value has no purpose but detecting errors and thus is |NULL_RTX|.

Now the compiler knows the insn code of this special insn but it needs to define the builtin as well. For special builtin functions this is not done automatically but in |altivec_init_builtins|:
\begin{lstlisting}[language=C++,basicstyle=\ttfamily\scriptsize,keywordstyle=\color{red}] 
def_builtin ("__builtin_altivec_mtvscr", void_ftype_v4si, ALTIVEC_BUILTIN_MTVSCR);
\end{lstlisting}
This adds |__builtin_altivec_mtvscr| to the list of defined functions and also gives the argument and return types (|void_ftype_v4si|, everything before |ftype| is the return type everything after the arguments). In this case it is not obvious why this needs to be done but there do exist builtins that have different insn codes depending on the used modes thus the type of the input arguments helps distinguish these differences. In this case only |v4si| is chosen as mode because we will also have an overloaded builtin for this new builtin function.
This is done similar to earlier by adding entries in |rs6000-c.c| for each combination of arguments and return types:
\begin{lstlisting}[language=C++,basicstyle=\ttfamily\scriptsize,keywordstyle=\color{red}] 
 { ALTIVEC_BUILTIN_VEC_MTVSCR, ALTIVEC_BUILTIN_MTVSCR,
   RS6000_BTI_void, RS6000_BTI_V4SI, 0, 0 },
 { ALTIVEC_BUILTIN_VEC_MTVSCR, ALTIVEC_BUILTIN_MTVSCR,
   RS6000_BTI_void, RS6000_BTI_unsigned_V4SI, 0, 0 },
 { ALTIVEC_BUILTIN_VEC_MTVSCR, ALTIVEC_BUILTIN_MTVSCR,
   RS6000_BTI_void, RS6000_BTI_bool_V4SI, 0, 0 },
 { ALTIVEC_BUILTIN_VEC_MTVSCR, ALTIVEC_BUILTIN_MTVSCR,
   RS6000_BTI_void, RS6000_BTI_V8HI, 0, 0 },
 { ALTIVEC_BUILTIN_VEC_MTVSCR, ALTIVEC_BUILTIN_MTVSCR,
   RS6000_BTI_void, RS6000_BTI_unsigned_V8HI, 0, 0 },
 { ALTIVEC_BUILTIN_VEC_MTVSCR, ALTIVEC_BUILTIN_MTVSCR,
   RS6000_BTI_void, RS6000_BTI_bool_V8HI, 0, 0 },
 { ALTIVEC_BUILTIN_VEC_MTVSCR, ALTIVEC_BUILTIN_MTVSCR,
   RS6000_BTI_void, RS6000_BTI_pixel_V8HI, 0, 0 },
 { ALTIVEC_BUILTIN_VEC_MTVSCR, ALTIVEC_BUILTIN_MTVSCR,
   RS6000_BTI_void, RS6000_BTI_V16QI, 0, 0 },
 { ALTIVEC_BUILTIN_VEC_MTVSCR, ALTIVEC_BUILTIN_MTVSCR,
   RS6000_BTI_void, RS6000_BTI_unsigned_V16QI, 0, 0 },
 { ALTIVEC_BUILTIN_VEC_MTVSCR, ALTIVEC_BUILTIN_MTVSCR,
   RS6000_BTI_void, RS6000_BTI_bool_V16QI, 0, 0 }
\end{lstlisting}
The return type obviously should be the same for all entries since there is no return type thus |RS6000_BTI_void| as first entry. Next there are 3 modes with different submodes, because all integer vector modes are allowed as first argument. A normal mode means that the elements are signed integers and an unsigned mode has unsigned elements. |Bool| elements have a single |bool| variable at each element and pixel is used for graphic usage of the AltiVec extension. The |0| at the end marks the second and third argument unused.
This sets the last step to completing our special builtin function that has a normal overloaded part.

At last we define a shorter function name in |s2pp.h|:
\begin{lstlisting}[language=C++,basicstyle=\ttfamily\scriptsize,keywordstyle=\color{red}] 
#define vec_mtvscr __builtin_vec_mtvscr
\end{lstlisting}
This completes our special intrinsic function!

For implementing a builtin function for earlier GCC versions (4.4 and earlier) I highly recommend the guide by Mauricio Alvarez \todo{(http://people.ac.upc.edu/ alvarez/media/ gcc-isa-extensions.html)} 

% chapter Y

% discussion
\chapter{Discussion}
\label{chapter:discussion}

Discussion...


% outlook
\chapter{Outlook}
\label{chapter:outlook}

As it was hinted in the discussion of this report, there are plans on extending the rs/6000 back-end of GCC so it supports the PPU's vector unit
This is surrently in the works currently as there are even some PPU-specific intrinsics already. Because the PPU is a simple PowerISA chip with a custom vector unit the implementation of intrinsics is easily transferable from the AltiVec vector extension to the PPU by swapping the |ALTIVEC| keyword with |S2PP|. With some knowledge of RTL this report then allows to add new composite intrinsic functions based on the instruction set of the PPU. By then this guide should be accompanied by a more complex thesis that explains the extension of the rs/6000 back-end and is meant to provide the reader with more insight into things like RTL and maybe makes reading the GCC internal obsolete.

%%%%%%%%%%%%%%%%%%%%%%%%%%%%%%%%%%%%%%%%





\appendix

\addcontentsline{toc}{chapter}{Appendix}
\chapter*{Appendix}

\section*{Acronyms}
\begin{acronym}
    \acro{ADC}{Analogue Digital Converter}
    \acro{ALU}{Arithmetic Logic Unit}
    \acro{ASIP}{Application Specific Instruction set Processor}
    \acro{asm}{Assembly}
    \acro{CADC}{Correlation Analogue Digital Converter}
    \acro{CISC}{Complex Instruction Set Computer}
    \acro{CPU}{Central Processing Unit}
    \acro{CR}{Conditional Register}
    \acro{DLS}{Digital Learning System}
    \acro{DRAM}{Dynamic RAM}
    \acro{insn}{instruction}
    \acro{IR}{Intermediate Representation}
    \acro{FPGA}{Field Programmable Gate Array}
    \acro{FPR}{Floating Point Register}
    \acro{FPU}{Ploating Point Unit}
    \acro{GCC}{GNU Compiler Collection}
    \acro{GPP}{General Purpose Processor}
    \acro{GPR}{General Purpose Register}
    \acro{HICANN}{High Input Count Neural Network}
    \acro{HICANN-DLS}{High Input Count Neural Network - Digital Learning System}
    \acro{ISA}{Instruction Set Architecture}
    \acro{LLVM}{Low Level Virtual Machine}
    \acro{LR}{Linker Register}
    \acro{LRA}{Local Register Allocator}
    \acro{MMU}{Memory Management Unit}
    \acro{MSB}{Most Significant Bit}
    \acro{nux}{alternative name for PPU}
    \acro{POWER}{Performance Optimization With Enhanced RISC}
    \acro{PPU}{Plasiticty Processing Unit}
    \acro{RAM}{Random Access Memory}
    \acro{RF}{Register File}
    \acro{RTL}{Register Transfer Language}
    \acro{RISC}{Reduced Instruction Set Computer}
    \acro{rs6000}{RISC system/6000}
    \acro{s2pp}{synaptic plasiticity processor}
    \acro{SIMD}{Single Input Multiple Data}
    \acro{SPR}{Special Purpose Register}
    \acro{SRAM}{Static RAM}
    \acro{VE}{Vector Extension}
    \acro{VR}{Vector Register}
    \acro{VSCR}{Vector Status and Control Register}
    \acro{VRSAVE}{Vector Save/Restore register}
    \acro{VRF}{Vector Register File}
    \acro{VMX}{Vector Media eXtension}

\end{acronym}

\section*{List of nux Intrinsics}

\begin{table}
    \adjustbox{max width=\textwidth}{%
        \begin{tabular}{c c¡c¡c¡c¡c¡p{5cm}}
            \multicolumn{1}{c}{intrinsic name} & \multicolumn{1}{c}{use} & \multicolumn{4}{c}{data types} & \multicolumn{1}{c}{effect}\\
            \cline{3-6}
             & & d & a & b & c & \\
            \hline \hline
                \begin{tabular}[x]{@{}c@{}}|fxv_add|\\|vec_add|\end{tabular} & |d = fxv_add(a,b)| & same as |a| & 
                \begin{tabular}[x]{@{}c@{}} |vector signed char|\\
                                            |vector unsigned char|\\
                                            |vector signed short|\\
                                            |vector unsigned short|\end{tabular}
                                            & same as |a| & &  add |a| and |b| modulo element-wise and write the result in |d|\\ 
            \cline{3-6}
                \begin{tabular}[x]{@{}c@{}}|fxv_sub|\\|vec_sub|\end{tabular} & |d = fxv_sub(a,b)| & same as |a| & 
                \begin{tabular}[x]{@{}c@{}} |vector signed char|\\
                                            |vector unsigned char|\\
                                            |vector signed short|\\
                                            |vector unsigned short|\end{tabular}
                                            & same as |a| & &  subtract |b| from |a| modulo element-wise and write the result in |d|\\ 
            \cline{3-6}
                \begin{tabular}[x]{@{}c@{}}|fxv_mul|\\|vec_mul|\end{tabular} & |d = fxv_mul(a,b)| & same as |a| & 
                \begin{tabular}[x]{@{}c@{}} |vector signed char|\\
                                            |vector unsigned char|\\
                                            |vector signed short|\\
                                            |vector unsigned short|\end{tabular}
                                            & same as |a| & &  multiply |a| and |b| modulo element-wise and write the result in |d|\\ 
            \cline{3-6}
                \begin{tabular}[x]{@{}c@{}}|fxv_addfs|\end{tabular} & |d = fxv_addfs(a,b)| & same as |a| & 
                \begin{tabular}[x]{@{}c@{}} |vector signed char|\\
                                            |vector unsigned char|\\
                                            |vector signed short|\\
                                            |vector unsigned short|\end{tabular}
                                            & same as |a| & &  add |a| and |b| saturational element-wise and write the result in |d|\\ 
            \cline{3-6}
                \begin{tabular}[x]{@{}c@{}}|fxv_subfs|\end{tabular} & |d = fxv_subfs(a,b)| & same as |a| & 
                \begin{tabular}[x]{@{}c@{}} |vector signed char|\\
                                            |vector unsigned char|\\
                                            |vector signed short|\\
                                            |vector unsigned short|\end{tabular}
                                            & same as |a| & &  subtract |b| from |a| saturational element-wise and write the result in |d|\\ 
            \cline{3-6}
                \begin{tabular}[x]{@{}c@{}}|fxv_mulfs|\end{tabular} & |d = fxv_mulfs(a,b)| & same as |a| & 
                \begin{tabular}[x]{@{}c@{}} |vector signed char|\\
                                            |vector unsigned char|\\
                                            |vector signed short|\\
                                            |vector unsigned short|\end{tabular}
                                            & same as |a| & &  multiply |a| and |b| saturational element-wise and write the result in |d|\\ 
            \cline{3-6}
                \begin{tabular}[x]{@{}c@{}}|fxv_stax|\\|vec_st|\end{tabular} & |fxv_stax(a,b,c)| & & 
                \begin{tabular}[x]{@{}c@{}} |vector signed char|\\\\
                                            |vector unsigned char|\\\\
                                            |vector signed short|\\\\
                                            |vector unsigned short|\\\end{tabular}
                                            & |int| &
                \begin{tabular}[x]{@{}c@{}} |vector signed char*|\\
                                            |signed char*|\\
                                            |vector unsigned char*|\\
                                            |unsigned char*|\\
                                            |vector signed short*|\\
                                            |signed short*|\\
                                            |vector unsigned short*|\\
                                            |unsigned short|\end{tabular}
                                            &  |a| is stored to memory address |c + b|\\ 
            \cline{3-6}
                \begin{tabular}[x]{@{}c@{}}|fxv_outx|\end{tabular} & |fxv_outx(a,b,c)| & & 
                \begin{tabular}[x]{@{}c@{}} |vector signed char|\\\\
                                            |vector unsigned char|\\\\
                                            |vector signed short|\\\\
                                            |vector unsigned short|\\\end{tabular}
                                            & |int| &
                \begin{tabular}[x]{@{}c@{}} |vector signed char*|\\
                                            |signed char*|\\
                                            |vector unsigned char*|\\
                                            |unsigned char*|\\
                                            |vector signed short*|\\
                                            |signed short*|\\
                                            |vector unsigned short*|\\
                                            |unsigned short|\end{tabular}
                                            &  |a| is stored to synaptic address |c + b|\\ 
            \cline{3-6}
                \begin{tabular}[x]{@{}c@{}}|fxv_lax|\\|vec_ld|\end{tabular} & |d = fxv_lax(a,b)| & 
                \begin{tabular}[x]{@{}c@{}} |vector signed char|\\\\
                                            |vector unsigned char|\\\\
                                            |vector signed short|\\\\
                                            |vector unsigned short|\\\end{tabular}
                                            & |int| &
                \begin{tabular}[x]{@{}c@{}} |vector signed char*|\\
                                            |signed char*|\\
                                            |vector unsigned char*|\\
                                            |unsigned char*|\\
                                            |vector signed short*|\\
                                            |signed short*|\\
                                            |vector unsigned short*|\\
                                            |unsigned short|\end{tabular}
                                            & & |d| is read from memory address |a + b|\\ 
            \cline{3-6}
                \begin{tabular}[x]{@{}c@{}}|fxv_inx|\end{tabular} & |d = fxv_inx(a,b)| & 
                \begin{tabular}[x]{@{}c@{}} |vector signed char|\\\\
                                            |vector unsigned char|\\\\
                                            |vector signed short|\\\\
                                            |vector unsigned short|\\\end{tabular}
                                            & |int| &
                \begin{tabular}[x]{@{}c@{}} |vector signed char*|\\
                                            |signed char*|\\
                                            |vector unsigned char*|\\
                                            |unsigned char*|\\
                                            |vector signed short*|\\
                                            |signed short*|\\
                                            |vector unsigned short*|\\
                                            |unsigned short|\end{tabular}
                                            & & |d| is read from synaptic address |a + b|\\ 
            \cline{3-6}
                \begin{tabular}[x]{@{}c@{}}|fxv_sel|\end{tabular} & |d = fxv_sel(a,b,c)| & same as |a| & 
                \begin{tabular}[x]{@{}c@{}} |vector signed char|\\
                                            |vector unsigned char|\\
                                            |vector signed short|\\
                                            |vector unsigned short|\end{tabular}
                                            & same as |a| & |2-bit int| & select element from |a| if |c| applies for that index otherwise select |b|, store the result in |d|\\ 
            \cline{3-6}
                \begin{tabular}[x]{@{}c@{}}|vec_extract|\end{tabular} & |d = vec_exctract(a,b)| & 
                \begin{tabular}[x]{@{}c@{}} |signed char|\\
                                            |unsigned char|\\
                                            |signed short|\\
                                            |unsigned short|\end{tabular}
                                            &
                \begin{tabular}[x]{@{}c@{}} |vector signed char|\\
                                            |vector unsigned char|\\
                                            |vector signed short|\\
                                            |vector unsigned short|\end{tabular}
                                            & |int| & & |d| is read from synaptic address |a + b|\\ 





            \cline{3-6}
                \begin{tabular}[x]{@{}c@{}}|vec_insert|\end{tabular} & |d = vec_insert(a,b)| & 
                \begin{tabular}[x]{@{}c@{}} |signed char|\\
                                            |unsigned char|\\
                                            |signed short|\\
                                            |unsigned short|\end{tabular}
                                            &
                \begin{tabular}[x]{@{}c@{}} |vector signed char|\\
                                            |vector unsigned char|\\
                                            |vector signed short|\\
                                            |vector unsigned short|\end{tabular}
                                            & |int| & & |d| is read from synaptic address |a + b|\\ 
            \cline{3-6}
                \begin{tabular}[x]{@{}c@{}}|vec_promote|\end{tabular} & |d = vec_exctract(a,b)| & 
                \begin{tabular}[x]{@{}c@{}} |signed char|\\
                                            |unsigned char|\\
                                            |signed short|\\
                                            |unsigned short|\end{tabular}
                                            &
                \begin{tabular}[x]{@{}c@{}} |vector signed char|\\
                                            |vector unsigned char|\\
                                            |vector signed short|\\
                                            |vector unsigned short|\end{tabular}
                                            & |int| & & |d| is read from synaptic address |a + b|\\ 
            \cline{3-6}
                \begin{tabular}[x]{@{}c@{}}|vec_sh| \\ |fxv_sh|\end{tabular} & |d = fxv_sh(a,b)| & 
                \begin{tabular}[x]{@{}c@{}} |signed char|\\
                                            |unsigned char|\\
                                            |signed short|\\
                                            |unsigned short|\end{tabular}
                                            &
                \begin{tabular}[x]{@{}c@{}} |vector signed char|\\
                                            |vector unsigned char|\\
                                            |vector signed short|\\
                                            |vector unsigned short|\end{tabular}
                                            & |int| & & |d| is read from synaptic address |a + b|\\ 
            \cline{3-6}
                \begin{tabular}[x]{@{}c@{}}|vec_splat_s16| \\ |vec_splat_u16 \\ fxv_splatb|\end{tabular} & |d = fxv_splatb(a)| & 
                \begin{tabular}[x]{@{}c@{}} |signed char|\\
                                            |unsigned char|\\
                                            |signed short|\\
                                            |unsigned short|\end{tabular}
                                            &
                \begin{tabular}[x]{@{}c@{}} |vector signed char|\\
                                            |vector unsigned char|\\
                                            |vector signed short|\\
                                            |vector unsigned short|\end{tabular}
                                            & |int| & & |d| is read from synaptic address |a + b|\\ 
            \cline{3-6}
                \begin{tabular}[x]{@{}c@{}}|vec_splat_s8| \\ |vec_splat_u8 \\ fxv_splath|\end{tabular} & |d = fxv_splath(a)| & 
                \begin{tabular}[x]{@{}c@{}} |signed char|\\
                                            |unsigned char|\\
                                            |signed short|\\
                                            |unsigned short|\end{tabular}
                                            &
                \begin{tabular}[x]{@{}c@{}} |vector signed char|\\
                                            |vector unsigned char|\\
                                            |vector signed short|\\
                                            |vector unsigned short|\end{tabular}
                                            & |int| & & |d| is read from synaptic address |a + b|\\ 
            \cline{3-6}
                \begin{tabular}[x]{@{}c@{}}|fxv_cmp|\end{tabular} & |fxv_cmp(a)| & 
                \begin{tabular}[x]{@{}c@{}} |signed char|\\
                                            |unsigned char|\\
                                            |signed short|\\
                                            |unsigned short|\end{tabular}
                                            &
                \begin{tabular}[x]{@{}c@{}} |vector signed char|\\
                                            |vector unsigned char|\\
                                            |vector signed short|\\
                                            |vector unsigned short|\end{tabular}
                                            & |int| & & |d| is read from synaptic address |a + b|\\ 
            \cline{3-6}
                \begin{tabular}[x]{@{}c@{}}|fxv_mtac| \\ |fxv_mtacfs|\end{tabular} & |fxv_mtac(a)| & 
                \begin{tabular}[x]{@{}c@{}} |signed char|\\
                                            |unsigned char|\\
                                            |signed short|\\
                                            |unsigned short|\end{tabular}
                                            &
                \begin{tabular}[x]{@{}c@{}} |vector signed char|\\
                                            |vector unsigned char|\\
                                            |vector signed short|\\
                                            |vector unsigned short|\end{tabular}
                                            & |int| & & |d| is read from synaptic address |a + b|\\ 
            \cline{3-6}
                \begin{tabular}[x]{@{}c@{}}|fxv_addactacm| \\ |fxv_addactacf|\end{tabular} & |fxv_addactac(a)| & 
                \begin{tabular}[x]{@{}c@{}} |signed char|\\
                                            |unsigned char|\\
                                            |signed short|\\
                                            |unsigned short|\end{tabular}
                                            &
                \begin{tabular}[x]{@{}c@{}} |vector signed char|\\
                                            |vector unsigned char|\\
                                            |vector signed short|\\
                                            |vector unsigned short|\end{tabular}
                                            & |int| & & |d| is read from synaptic address |a + b|\\ 
            \cline{3-6}
                \begin{tabular}[x]{@{}c@{}}|fxv_addacm| \\ |fxv_addacfs|\end{tabular} & |fxv_addac(a)| & 
                \begin{tabular}[x]{@{}c@{}} |signed char|\\
                                            |unsigned char|\\
                                            |signed short|\\
                                            |unsigned short|\end{tabular}
                                            &
                \begin{tabular}[x]{@{}c@{}} |vector signed char|\\
                                            |vector unsigned char|\\
                                            |vector signed short|\\
                                            |vector unsigned short|\end{tabular}
                                            & |int| & & |d| is read from synaptic address |a + b|\\ 
            \cline{3-6}
                \begin{tabular}[x]{@{}c@{}}|fxv_mam| \\ |fxv_mafs|\end{tabular} & |fxv_ma(a)| & 
                \begin{tabular}[x]{@{}c@{}} |signed char|\\
                                            |unsigned char|\\
                                            |signed short|\\
                                            |unsigned short|\end{tabular}
                                            &
                \begin{tabular}[x]{@{}c@{}} |vector signed char|\\
                                            |vector unsigned char|\\
                                            |vector signed short|\\
                                            |vector unsigned short|\end{tabular}
                                            & |int| & & |d| is read from synaptic address |a + b|\\ 
            \cline{3-6}
                \begin{tabular}[x]{@{}c@{}}|fxv_matacm| \\ |fxv_matacfs|\end{tabular} & |fxv_matac(a)| & 
                \begin{tabular}[x]{@{}c@{}} |signed char|\\
                                            |unsigned char|\\
                                            |signed short|\\
                                            |unsigned short|\end{tabular}
                                            &
                \begin{tabular}[x]{@{}c@{}} |vector signed char|\\
                                            |vector unsigned char|\\
                                            |vector signed short|\\
                                            |vector unsigned short|\end{tabular}
                                            & |int| & & |d| is read from synaptic address |a + b|\\ 
            \cline{3-6}
                \begin{tabular}[x]{@{}c@{}}|fxv_multacm| \\ |fxv_multacfs|\end{tabular} & |fxv_multac(a)| & 
                \begin{tabular}[x]{@{}c@{}} |signed char|\\
                                            |unsigned char|\\
                                            |signed short|\\
                                            |unsigned short|\end{tabular}
                                            &
                \begin{tabular}[x]{@{}c@{}} |vector signed char|\\
                                            |vector unsigned char|\\
                                            |vector signed short|\\
                                            |vector unsigned short|\end{tabular}
                                            & |int| & & |d| is read from synaptic address |a + b|\\ 
            \cline{3-6}
                \begin{tabular}[x]{@{}c@{}}|fxv_addtac|\end{tabular} & |fxv_addtac(a)| & 
                \begin{tabular}[x]{@{}c@{}} |signed char|\\
                                            |unsigned char|\\
                                            |signed short|\\
                                            |unsigned short|\end{tabular}
                                            &
                \begin{tabular}[x]{@{}c@{}} |vector signed char|\\
                                            |vector unsigned char|\\
                                            |vector signed short|\\
                                            |vector unsigned short|\end{tabular}
                                            & |int| & & |d| is read from synaptic address |a + b|\\ 
            \cline{3-6}
                \begin{tabular}[x]{@{}c@{}}|fxv_pckbu| \\ |fxv_pckbl|\end{tabular} & |d = fxv_pckbu(a,b)| & 
                \begin{tabular}[x]{@{}c@{}} |signed char|\\
                                            |unsigned char|\\
                                            |signed short|\\
                                            |unsigned short|\end{tabular}
                                            &
                \begin{tabular}[x]{@{}c@{}} |vector signed char|\\
                                            |vector unsigned char|\\
                                            |vector signed short|\\
                                            |vector unsigned short|\end{tabular}
                                            & |int| & & |d| is read from synaptic address |a + b|\\ 
            \cline{3-6}
                \begin{tabular}[x]{@{}c@{}}|fxv_upckbl| \\ |fxv_upckbr|\end{tabular} & |d = fxv_upckbl(a,b)| & 
                \begin{tabular}[x]{@{}c@{}} |signed char|\\
                                            |unsigned char|\\
                                            |signed short|\\
                                            |unsigned short|\end{tabular}
                                            &
                \begin{tabular}[x]{@{}c@{}} |vector signed char|\\
                                            |vector unsigned char|\\
                                            |vector signed short|\\
                                            |vector unsigned short|\end{tabular}
                                            & |int| & & |d| is read from synaptic address |a + b|\\ 

        \end{tabular}
    }
    \caption{\label{appendix:builtinlist} List of all implemented built-ins and how they are used.}
\end{table}


\begin{tabular}{l l p{9cm}}
    mnemonic & operands & description \\
    \hline
    |add| & |RT, RA, RB| & \textbf{add} |RB| to |RA| and store the result in |RT| \\
    |addi| & |RT, RA, SI| & \textbf{add} |SI| to |RA| and store the result in |RT| \\
    |addis| & |RT, RA, SI| & \textbf{add} |SI| \textbf{s}hifted left by 16 bit to |RA| and store the result in |RT| \\
    |and| & |RA, RS, RB| & |RS| and |RB| are \textbf{and}ed and the result is stored in |RT| \\
    |b| & |target_addr| & \textbf{b}ranch to the code at |target_addr|\\
    |ble| & |BF, target_addr| & \textbf{b}ranch to the code at |target_addr| if |BF| is \textbf{l}ess or \textbf{e}qual \\
    |blr| & & \textbf{b}ranch to the code at address in the \textbf{l}inker \textbf{r}egister \\
    |cmp| & |BF, L, RA, RB| & |RA| and |RB| are \textbf{c}o\textbf{mp}ared and the result (|gt,lt,eq|) is stored in |BF|, |L| depicts if 32-bit or 64-bit are compared \\
    |cmplwi| & |BF, RA, SI| & |RA| \textbf{c}o\textbf{mp}ared \textbf{l}ogically \textbf{w}ordwise with \textbf{i}mmediate |SI| and the result is stored in |BF|\\
    |and| & |RA, RS, RB| & |RS| and |RB| are \textbf{and}ed and the result is stored in |RT| \\
    |eieio| & & enforce in-order execution of I/O\\
    |isync| & & instruction cache synchronize\\
    |la| & |RT, D(RA)| & \textbf{l}oad \textbf{a}ggregate |D + RA| into |RT|\\
    |li| & |RT, SI| & \textbf{l}oad \textbf{i}mmediate value |SI| into |RT|\\
    |lis| & |RT, SI| & \textbf{l}oad \textbf{i}mmediate value |SI| \textbf{s}hifted left by 16 bit into |RT|\\
    |lbz| & |RT, D(RA)| & \textbf{l}oad \textbf{b}yte at address |D+RA| into |RT|, fill the other bits with \textbf{z}eros \\
    |lwz| & |RT, D(RA)| & \textbf{l}oad \textbf{w}ord at address |D+RA| into |RT|, fill the other bits with \textbf{z}eros \\
    |mflr| & |RT| & \textbf{m}ove \textbf{f}rom \textbf{l}inker \textbf{r}egister to |RT|\\
    |mr| & |RT, RA| & \textbf{m}ove \textbf{r}egister |RA| to |RT| \\
    |nop| & & halts execution until the previous instructions are finished EDITHIS\\
    |rlwinm| & |RA, RS, SH, MB, ME| & \textbf{r}otate \textbf{l}eft \textbf{w}ord in |RS| by \textbf{i}mmediate |SH| bits then a\textbf{n}d with \textbf{m}ask which is 1 from |MB+32| to |ME+32| and 0 else, store to |RA|\\
    |stw| & |RS, D(RA)| & \textbf{st}ore \textbf{w}ord from |RS| to address |D+RA|\\
    |stwu| & |RS, D(RA)| & \textbf{st}ore \textbf{w}ord from |RS| to address |D+RA| and \textbf{u}pdate RA to |D+RA|\\
    |sync| & & synchronize data cache\\
\end{tabular}
\add{caption and reference to PPC book and asm website correct the table}

Obviously the mnemonics follow a certain pattern that has letters which can be interchanged to alter the meaning of the mnemonic, some of these characters are:
\begin{description}
    \lstitem{i} indicates that the instructions uses an immediate value
    \lstitem{b} stands for byte and references the size of the operand
    \lstitem{h} stands for halfword and references the size of the operand
    \lstitem{w} stands for word and references the size of the operand
    \lstitem{s} indicates that one of the operands is shifted
    \lstitem{g, ge, l, le, e} stand for greater, greater or equal, less, less or equal and equal which is the possible content of the conditional register
\end{description}

There are also special operands which might occur in |asm| which behave like pointers:
\begin{description}
    \lstitem{@l(C)} is equivalent to the lower order 16 bits of of |C| in the symbol table
    \lstitem{@ha(C)} is equivalent to the higher order 16 bits of of |C| in the symbol table and minds the sign extension
\end{description}

Additionally there exist markers which are intended for debugging:
\begin{description}
        \lstitem{.loc # # #} marks a line of code (file, line, column) in the source file
        \lstitem{.LVL} is a local label which can be discarded
        \lstitem{.LFB} marks the begin of a function
        \lstitem{.LFE} marks the end of a function
        \lstitem{.LC0} is a constant of the literals table at position |0|
\end{description}

\add{all these tables in the appendix}
\newpage
\listoftodos[Notes]



\pagestyle{plain}       % ab hier werden nur noch seitenzahlen angezeigt
% ------------- Abkuerzungsverzeichnis -----------
%\cleardoublepage
%\addcontentsline{toc}{chapter}{Nomenclature}
\printnomenclature
% ------------- Abbildungsverzeichnis ------------
%\cleardoublepage
%\addcontentsline{toc}{chapter}{List Of Figures}
%\listoffigures
% -------------- Tabellenverzeichnis -------------
%\cleardoublepage
%\addcontentsline{toc}{chapter}{List Of Tables}
%\listoftables
% ------------- Algorithmenverzeichnis -----------
%\cleardoublepage
%\addcontentsline{toc}{chapter}{List Of Algorithms}
%\listofalgorithms
% ------------------------------------------------

\pagestyle{headings}    % ab hier werden seitenzahlen UND kapitelueberschriften in der kopfzeile angezeigt

% ------------- Bibliography ---------------------
\cleardoublepage
\bibliographystyle{bib/agu}
%\bibliography{bib/vision}
\bibliography{bib/myrefs}
\addcontentsline{toc}{chapter}{Bibliography}
% ------------------------------------------------


\pagestyle{plain}

% ---------------------- Index -------------------
%\cleardoublepage
%\addcontentsline{toc}{chapter}{Index}
\printindex
% ------------------------------------------------

% ------------- Thanks ... -----------------------

%\cleardoublepage
%\addcontentsline{toc}{chapter}{Acknowledgments}
\include{tex/acknowledgments}

% ------------------------------------------------


% ------------- Statement ------------------------

\cleardoublepage
\pagestyle{empty}
\selectlanguage{ngerman}
\thispagestyle{empty}
\parindent0pt
\vspace*{1cm}

{\Large{\bf{Statement of Originality (Erkl"arung):}}}
\vspace{1cm}
\selectlanguage{english}
\\
I certify that this thesis, and the research to which it refers, are the product of my own work. 
Any ideas or quotations from the work of other people, published or otherwise, 
are fully acknowledged in accordance with the standard referencing practices of the discipline. 
%I certify that this thesis has not been submitted for a higher degree to any other university or institution. 

\vspace{2cm}
\selectlanguage{ngerman}
Ich versichere, dass ich diese Arbeit selbst"andig verfasst und keine anderen als die
angegebenen Quellen und Hilfsmittel benutzt habe.

\vspace{3cm}
\selectlanguage{english}

Heidelberg, \today \\
\hspace*{10.0cm}.......................................\\
\hspace*{11.1cm}(signature)


% ------------------------------------------------




\end{document}
